\section{\mgp: A Decision-Level Diagnostic}
\label{sec:mgp6}

\mgp{} is an author-designed diagnostic slice for write-time memory governance.
We use the word \emph{diagnostic} deliberately and throughout: it is a pilot
instrument for probing decision mechanisms, not a community benchmark, an
independent test collection, or a measure of general QA or retrieval efficiency. All scenarios are synthetic,
English-only, and low-risk, and contain no real personal data.

\subsection{v4.3 final release and leakage controls}

The historical v2 pilot remains in the artifact as a record of the earlier
instrument and its limitations. The successor release, \mgp{} v4.3.0, is a
human-authored, template-guided, controlled synthetic diagnostic for the three
decision actions \keep{}, \supersede{}, and \abstain{}. It contains 900 items
in 100 semantic components across 20 domains. Components, rather than
individual items, are assigned intact to train (180 items; 20 components),
development (180; 20), and hidden (540; 60) partitions. Each split is exactly
balanced across the three actions, and each domain contributes one component to
train, one to development, and three to hidden.

This redesign follows the v2 shortcut diagnosis: earlier rewrite-oriented
versions retained label-correlated lexical and structural cues, so held-out item
performance could reflect scenario recognition rather than the intended
reasoning about authority, scope, time, identity, and prior memory. The v4.3
content architecture uses balanced factual evidence and controlled distractors.
It passed duplicate, near-duplicate, model-field allowlist, and shortcut-audit
gates; the detailed restricted-view and shallow-feature audits remain private.
Public model-facing records exclude gold labels, design targets, private IDs,
component metadata, and annotator metadata.

Two completed human annotation packets cover all 900 items. They agree on all
actions (raw agreement $=1.0$, Cohen's $\kappa=1.0$), with no exclusions or
label adjudications. The 300 \abstain{} items are intentionally ambiguous
scenarios, not gold-label disagreements. Completion is accepted through the
dataset-owner attestation: the release builder verifies records, coverage,
hashes, and agreement, but does not independently authenticate reviewer identity
or annotation duration.

\subsection{Conflict taxonomy and distribution}

The pilot contains $180$ scenarios, $30$ for each of six scenario families,
with an action imbalance of
$90$ \supersede, $66$ \abstain, and $24$ \keep{} that we handle by reporting
macro-F1 and per-action results rather than accuracy alone.

\begin{center}
\small
\begin{tabular}{@{}llr l@{}}
\toprule
Type & Diagnostic condition & Rows & Gold action(s) \\
\midrule
T1 & preference drift over time & 30 & 30 \supersede \\
T2 & episodic factual correction & 30 & 30 \supersede \\
T3 & stable-preference replacement & 30 & 30 \supersede \\
T4 & ambiguity / referent uncertainty & 30 & 30 \abstain \\
T5 & transient capture trap & 30 & 24 \keep, 6 \abstain \\
T6 & confidence laundering & 30 & 30 \abstain \\
\bottomrule
\end{tabular}
\end{center}

\noindent
T1--T3 target supersession under different durability regimes; a subset of $10$
rows per type is \emph{cue-free} (\textsf{implicit\_conflict}), i.e.\ the
perturbation contains none of the explicit change cues (``now,'' ``instead of,''
``actually,'' \dots) yet still fixes a mutually exclusive current value. T4/T6
require \abstain{} because evidence is referentially uncertain or laundered
through an unreliable source; T5 separates a durable preference (\keep) from a
temporary override. Some supersession cases also have minimal-pair T4 partners.
These pairs test a narrow cue change, but they do not remove the global shortcut:
conflict type still determines the action throughout v2.

\subsection{Query modes}

Every scenario carries one query mode used only at the answer stage:
\emph{direct recall} ($108$ items), \emph{premise resistance} ($45$; the query
embeds the stale/unsupported value as a false presupposition that a correct
answer must reject), and \emph{implicit application} ($27$; a concrete task
solvable only by using the active value, with neither value named in the query).
These modes let us measure whether a system that makes the right write decision
also refuses stale premises and applies the active value in context.

\subsection{Splits}

Splits are $60$ train / $30$ dev / $90$ test, stratified so each type contributes
$10/5/15$ rows and T5 contributes $8/2$, $4/1$, $12/3$ \keep/\abstain. No
template-family identifier and no minimal-pair group crosses a split boundary,
which blocks literal family overlap. It does \emph{not} block construct leakage
when the same family-to-action rule holds in every split. We therefore treat the
$90$-item test split as a held-out item split only, never as independent
validation.

\subsection{Pre-specified validity checks}

Before comparing policies we fit lookup baselines on train+dev and evaluate them
on test using only (i) the six-way type, (ii) the finer conflict type, and (iii)
source/reliability/temporal metadata. We also compute macro-accuracy and paired
bootstrap intervals by \texttt{template\_family\_id}; for
\texttt{minimal\_pair\_group\_id}, unpaired items are singleton clusters. A
shortcut baseline that approaches the full policy is a failed construct-validity
check, not a competitive result.

\subsection{v4.3 evaluation contract}

For v4.3, systems submit exactly one action for every public identifier using
the two fields \texttt{public\_id} and \texttt{action}; the exact JSON form is
part of the released evaluation protocol. The evaluator rejects missing,
duplicate, unknown, and extra identifiers. It reports
accuracy, macro-F1, per-class precision/recall/F1, a confusion matrix, coverage,
and component-cluster bootstrap intervals. Items within a semantic component are
not treated as independent; paired comparisons use component-aware resampling
or component-level analysis. LongMemEval is reported separately as external
memory evaluation, not as a substitute for this controlled diagnostic.

\subsection{Schema, invariants, and scoring}

Each row provides a dialogue prefix, prior memory state, perturbation event,
source/temporal metadata, a gold action, an auditable gold memory state, the
query mode, the expected answer, and a regex judge specification with required
(\texttt{include}) and forbidden (\texttt{not\_include}) spans. Evaluation is
two-stage: the write stage sees only the prefix, prior state, and perturbation
and commits an action and predicted state \emph{before} the query is revealed;
the answer stage then answers from the predicted state. A validator enforces the
state-transition invariants of Section~\ref{sec:actionspace} (one active record
per slot, retained superseded predecessors, single-successor lineage,
byte-equivalent \abstain{} states, active-only answers), and a released corpus
prints \texttt{VALIDATION PASSED: 180 examples, 0 errors}.

We score two families of metric (full definitions in Appendix~\ref{app:metrics}):
\emph{decision-level} --- action accuracy, per-action and macro-F1, and the
action false-commit rate on gold-\abstain{} items --- and \emph{memory-state}
--- state invalid-commit on gold-\abstain{} items, five-condition lineage
integrity on \supersede{} items, and stale-retrieval rate (the fraction of answers surfacing a forbidden
superseded value). Two shipped forgetting-rate metrics saturate at $1.0$ for every
system and are excluded from all tables as uninformative.

\subsection{Annotation: one independent annotator plus a model audit panel}
\label{sec:annotation}

We are precise about the annotation evidence, which has two separate components
and must never be described as multi-annotator human agreement.

\begin{description}[leftmargin=0em,topsep=2pt,itemsep=2pt]
\item[One independent human annotator] re-labelled the $90$-item test split
against the annotation manual. Raw agreement with gold is $0.9889$
($\text{CI}_{95}\,[0.9667, 1.0]$), Cohen's $\kappa = 0.98146$
($\text{CI}_{95}\,[0.93997, 1.0]$), with exactly one disagreement
(\texttt{MGP6V2-T3-TEST-016}, an \textsf{implicit\_conflict} \supersede$\to$\abstain{}
case). This is $n{=}1$ annotator; the annotator's human status rests on author
attestation, not on evidence in the released files.
\item[A four-model audit panel] (four distinct LLMs) independently labelled the
same $90$ items: Fleiss $\kappa = 0.9815$ across the panel, with the majority
vote matching gold on all $90$ items ($88$ unanimous). This is a model audit, not
human inter-annotator agreement.
\end{description}
The larger human review sheet ($180$ rows) exists but is unfilled; a genuine
multi-annotator human release is future work (Section~\ref{sec:limitations}).
The panel and the single annotator may share the manual's reading because the
labelling policy is unambiguous by construction --- a construct-validity caveat we
carry explicitly.

\subsection{Construction grounding}

Scenario \emph{shapes} were informed by inspecting real corpora --- LongMemEval
knowledge-update items and its false-premise (\texttt{\_abs}) questions, Memora
operation chains, GateMem's machine-readable judge structure, and named
confidence-laundering patterns --- but every released row is a newly authored
synthetic paraphrase; no source content is reused. This grounds the failure modes
without importing personal data, and it is the reason \mgp{} shares construct
names (implicit conflict, premise resistance) with concurrent read-side work
\citep{chao2026stale}, which we cite to disambiguate.
